\documentclass[11pt]{article}
\usepackage{listings}
\usepackage{xcolor}
\usepackage{graphicx}
\usepackage{subcaption}
\usepackage{booktabs}
\usepackage{float}
\usepackage{titling}
\setlength{\droptitle}{-3cm}
\lstset{
    language=Python,
    basicstyle=\ttfamily\small,
    keywordstyle=\color{blue},
    commentstyle=\color{gray},
    stringstyle=\color{green!50!black},
    backgroundcolor=\color{gray!5},
    frame=single,
    rulecolor=\color{gray!40},
    breaklines=true,
    captionpos=b,
    showstringspaces=false
}

\title{\textbf{Labwork 3 - CUDA}}
\author{
Do Minh Quang - 2540095\\
University of Science and Technology of Hanoi
}
\date{}

\begin{document}
\maketitle

\section{RGB to Grayscale Conversion}
This labwork is about converting an RGB image to Grayscale image on CPU and GPU in order to observe the boost up of GPU over CPU. The experiment uses 2 RGB images of different resolutions, $256 \times 256$ and $3167 \times 3167$ to examine how image size affects execution time. 

\begin{figure}[H]
    \centering
    \includegraphics[width=0.5\textwidth]{figure/bird.jpg}
    \caption{RGB image.}
\end{figure}

The image is first loaded and reshaped into 2D array of shape $(N, 3)$, where $N$ is the total number of pixels.

\subsection{CPU Implementation}
The CPU implementation processes pixels sequentially in the \texttt{for} loop. For each pixel, the average of the RGB channels is calculated and stored in the output array.

\begin{lstlisting}[language=Python, label={lst:code}]
    def rgb2gray_cpu(rgb, gray):
        for i in range(rgb.shape[0]):
            g = np.uint8((int(rgb[i, 0]) + int(rgb[i, 1]) + int(rgb[i, 2])) / 3)
            gray[i, 0] = gray[i, 1] = gray[i, 2] = g
\end{lstlisting}

\subsection{GPU Implementation}
For the GPU version, the flattened RGB image is first transferred from CPU memory to the GPU memory. The CUDA kernel, \texttt{rgb2gray\_gpu}, is then launched to process pixels in parallel, with each thread handling one pixel. Finally the result grayscale image is copied back to CPU memory for visualization. The initial block size is set to 64 threads. 

\begin{lstlisting}[language=Python, label={lst:code}]
    def grayscale_gpu(flatten_rgb, blockSize):
        gridSize = (pixelCount + blockSize - 1) // blockSize
        rgb_gpu = cuda.to_device(flatten_rgb)
        gray_gpu = cuda.device_array((pixelCount, 3), np.uint8)
        rgb2gray_gpu[gridSize, blockSize](rgb_gpu, gray_gpu)
        host_gray_gpu = gray_gpu.copy_to_host()
        return host_gray_gpu
\end{lstlisting}

\section{Results}

\subsection{CPU and GPU performance}
For the quantitative results, Table~\ref{tab:performance} shows that the GPU achieved a speedup of approximately 2 times for the $256 \times 256$ image and 188 times for the $3167 \times 3167$ image. These results shows that the performance difference between CPU and GPU becomes noticeable only when processing sufficiently large images.

\begin{table}[H]
\centering
\caption{Execution time of CPU and GPU implementations for different image sizes.}
\label{tab:performance}
\begin{tabular}{lrr}
\toprule
\textbf{Image Size} & \textbf{CPU (s)} & \textbf{GPU (s)} \\
\midrule
$256 \times 256$ & 0.1841 & 0.0872 \\
$3167 \times 3167$ & 21.1528 & 0.1126 \\
\bottomrule
\end{tabular}
\end{table}

For the qualitative result, both implementations produced visually consistent grayscale images. However, the GPU image version appears slightly sharper, possibly due to differences in data types or numerical operations during computation, as shown in Figure~\ref{fig:qualitative}.

\begin{figure}[H]
\centering
\begin{subfigure}[t]{0.47\textwidth}
\centering
\includegraphics[width=\textwidth]{figure/gray_cpu.jpg}
\caption{CPU grayscale output.}
\end{subfigure}
\hfill
\begin{subfigure}[t]{0.47\textwidth}
\centering
\includegraphics[width=\textwidth]{figure/gray_gpu.jpg}
\caption{GPU grayscale output.}
\end{subfigure}
\caption{Comparison of grayscale images produced by the CPU and GPU.}
\label{fig:qualitative}
\end{figure}

\subsection{Impact of block size}
In this part the 3167x3167 image is used to evaluate the impact of block size on GPU execution time. The block size is varied from 1 to 1024 threads, and the corresponding execution times are recorded. The results are presented in Figure~\ref{fig:block_size}. 

\begin{figure}[H]
\centering
\includegraphics[width=0.7\textwidth]{figure/blocksize_vs_time.png}
\caption{GPU execution time for different block sizes.}
\label{fig:block_size}
\end{figure}

The measured running time generally decreases as the block size increases. However, the performance gain becomes marginal between block size of 128 and 1024. Among the tested configurations, a block size of 512 achieved the shortest execution time of 0.015299 seconds, while a block size of 1 resulted in the longest time of 0.095496 seconds.


\section{Conclusion}
The results show that GPU parallelism provides greater performance benefits for sufficiently large images. Block size also affects GPU performance, with the optimal configuration depending on the workload and hardware architecture.

\end{document}